\documentclass[9pt,t]{beamer}
% =====================================================================
% finance-beamer.tex — shared Beamer style for the LLM-in-Finance course
% =====================================================================
% Reusable slide style. Each deck \input's this immediately after
% \documentclass{beamer}. It provides:
%   * the Madrid theme + book colour palette (matches the printed book)
%   * two book-style framing blocks:
%       \begin{bigpicture} ... \end{bigpicture}   ("the bigger picture")
%       \begin{underhood}   ... \end{underhood}    ("under the hood")
%     These mirror the book's `context` and `deepdive` tcolorboxes so the
%     same intuition-vs-mechanism scaffold the reader meets in the book
%     reappears on the slides.
%   * a Python listings style for code frames
%   * appendix conventions: \appendixstart drops a divider and switches
%     to non-numbered backup frames so "extra / less central" material
%     lives after the main talk without inflating the slide count.
%
% Usage:
%   \documentclass{beamer}
%   % =====================================================================
% finance-beamer.tex — shared Beamer style for the LLM-in-Finance course
% =====================================================================
% Reusable slide style. Each deck \input's this immediately after
% \documentclass{beamer}. It provides:
%   * the Madrid theme + book colour palette (matches the printed book)
%   * two book-style framing blocks:
%       \begin{bigpicture} ... \end{bigpicture}   ("the bigger picture")
%       \begin{underhood}   ... \end{underhood}    ("under the hood")
%     These mirror the book's `context` and `deepdive` tcolorboxes so the
%     same intuition-vs-mechanism scaffold the reader meets in the book
%     reappears on the slides.
%   * a Python listings style for code frames
%   * appendix conventions: \appendixstart drops a divider and switches
%     to non-numbered backup frames so "extra / less central" material
%     lives after the main talk without inflating the slide count.
%
% Usage:
%   \documentclass{beamer}
%   % =====================================================================
% finance-beamer.tex — shared Beamer style for the LLM-in-Finance course
% =====================================================================
% Reusable slide style. Each deck \input's this immediately after
% \documentclass{beamer}. It provides:
%   * the Madrid theme + book colour palette (matches the printed book)
%   * two book-style framing blocks:
%       \begin{bigpicture} ... \end{bigpicture}   ("the bigger picture")
%       \begin{underhood}   ... \end{underhood}    ("under the hood")
%     These mirror the book's `context` and `deepdive` tcolorboxes so the
%     same intuition-vs-mechanism scaffold the reader meets in the book
%     reappears on the slides.
%   * a Python listings style for code frames
%   * appendix conventions: \appendixstart drops a divider and switches
%     to non-numbered backup frames so "extra / less central" material
%     lives after the main talk without inflating the slide count.
%
% Usage:
%   \documentclass{beamer}
%   \input{../_style/finance-beamer.tex}
%   \title{...}\subtitle{...}\author{...}\institute{...}\date{\today}
%   \begin{document} ... \end{document}
% =====================================================================

\usetheme{Madrid}
\usecolortheme{whale}
\setbeamertemplate{navigation symbols}{}
\setbeamertemplate{caption}[numbered]
\setbeamercovered{transparent}

% --- Density controls: keep dense, equation-heavy frames on one slide ---
% Slightly smaller body text + tighter lists/blocks. Frame titles and the
% Madrid chrome keep their normal size, so the deck still reads as Madrid,
% just with more vertical room for content.
\setbeamerfont{normal text}{size=\small}
\setbeamertemplate{itemize/enumerate body begin}{\small}
\setbeamertemplate{itemize/enumerate subbody begin}{\footnotesize}
% Tighter block spacing.
\setbeamertemplate{block begin}{%
  \par\vskip2pt%
  \begin{beamercolorbox}[colsep*=.75ex,leftskip=1ex,rightskip=1ex]{block title}%
    \usebeamerfont*{block title}\insertblocktitle%
  \end{beamercolorbox}%
  {\parskip0pt\vskip-1pt}%
  \begin{beamercolorbox}[colsep*=.75ex,vmode,leftskip=1ex,rightskip=1ex]{block body}%
  \vskip1pt}
\setbeamertemplate{block end}{\end{beamercolorbox}\vskip2pt}

\usepackage{amsmath, amssymb}
\usepackage{booktabs}
\usepackage{xcolor}
\usepackage{listings}
\usepackage[skins, breakable]{tcolorbox}

% --- Book colour palette (kept in sync with book/preamble.tex) ---
\definecolor{linkblue}{RGB}{14, 75, 140}
\definecolor{deepdivebg}{RGB}{235, 244, 255}
\definecolor{narrativebg}{RGB}{255, 252, 240}
\definecolor{narrativerule}{RGB}{180, 130, 40}
\definecolor{steelblue}{RGB}{70, 130, 180}
\definecolor{forestgreen}{RGB}{34, 139, 34}
\definecolor{crimson}{RGB}{220, 20, 60}
\definecolor{darkorange}{RGB}{255, 140, 0}
\definecolor{codebg}{RGB}{248, 248, 248}
\definecolor{coderule}{RGB}{210, 210, 210}
\definecolor{keywordblue}{RGB}{0, 100, 180}
\definecolor{commentgray}{RGB}{120, 120, 120}
\definecolor{stringgreen}{RGB}{30, 130, 30}

% Tie the Madrid structural colour to the book's link blue.
\setbeamercolor{structure}{fg=linkblue}
\setbeamercolor{title}{fg=white}
\setbeamercolor{frametitle}{fg=white}

% --- "the bigger picture" : narrative / intuition framing -----------
\newtcolorbox{bigpicture}{%
  enhanced, breakable, boxsep=2pt,
  colback  = narrativebg,
  colframe = narrativerule,
  boxrule  = 0pt, toprule = 1.4pt, bottomrule = 0.4pt,
  leftrule = 0pt, rightrule = 0pt, arc = 0pt,
  left = 6pt, right = 6pt, top = 4pt, bottom = 5pt,
  fonttitle = \footnotesize\scshape\color{narrativerule},
  title = {the bigger picture}, attach title to upper = {\par\smallskip},
}

% --- "under the hood" : the mechanism / technical detail ------------
\newtcolorbox{underhood}{%
  enhanced, breakable, boxsep=2pt,
  colback  = deepdivebg,
  colframe = linkblue,
  boxrule  = 0pt, toprule = 1.4pt, bottomrule = 0.4pt,
  leftrule = 0pt, rightrule = 0pt, arc = 0pt,
  left = 6pt, right = 6pt, top = 4pt, bottom = 5pt,
  fonttitle = \footnotesize\scshape\color{linkblue},
  title = {under the hood}, attach title to upper = {\par\smallskip},
}

% --- Python code style ---------------------------------------------
\lstset{
  basicstyle    = \ttfamily\footnotesize,
  keywordstyle  = \color{keywordblue}\bfseries,
  commentstyle  = \color{commentgray}\itshape,
  stringstyle   = \color{stringgreen},
  showstringspaces = false,
  language      = Python,
  breaklines    = true,
  backgroundcolor = \color{codebg},
  frame         = single,
  rulecolor     = \color{coderule},
  numbers       = none,
  aboveskip     = 4pt, belowskip = 4pt,
}

% --- Appendix conventions ------------------------------------------
% \appendixstart{Title} : begins the "extra material" appendix. Frames
% after it are not counted toward the main deck's frame total, keeping
% the core talk lean while preserving the deeper / optional content.
\newcommand{\appendixstart}[1]{%
  \appendix
  \setbeamertemplate{frametitle continuation}{}%
  \begin{frame}[noframenumbering]
    \begin{center}
      {\Large\scshape\color{linkblue} Appendix}\\[4pt]
      {\large #1}\\[10pt]
      {\footnotesize Extra / optional material — beyond the core lecture.}
    \end{center}
  \end{frame}
}
% Use \begin{frame}[noframenumbering]{...} for each appendix frame.

%   \title{...}\subtitle{...}\author{...}\institute{...}\date{\today}
%   \begin{document} ... \end{document}
% =====================================================================

\usetheme{Madrid}
\usecolortheme{whale}
\setbeamertemplate{navigation symbols}{}
\setbeamertemplate{caption}[numbered]
\setbeamercovered{transparent}

% --- Density controls: keep dense, equation-heavy frames on one slide ---
% Slightly smaller body text + tighter lists/blocks. Frame titles and the
% Madrid chrome keep their normal size, so the deck still reads as Madrid,
% just with more vertical room for content.
\setbeamerfont{normal text}{size=\small}
\setbeamertemplate{itemize/enumerate body begin}{\small}
\setbeamertemplate{itemize/enumerate subbody begin}{\footnotesize}
% Tighter block spacing.
\setbeamertemplate{block begin}{%
  \par\vskip2pt%
  \begin{beamercolorbox}[colsep*=.75ex,leftskip=1ex,rightskip=1ex]{block title}%
    \usebeamerfont*{block title}\insertblocktitle%
  \end{beamercolorbox}%
  {\parskip0pt\vskip-1pt}%
  \begin{beamercolorbox}[colsep*=.75ex,vmode,leftskip=1ex,rightskip=1ex]{block body}%
  \vskip1pt}
\setbeamertemplate{block end}{\end{beamercolorbox}\vskip2pt}

\usepackage{amsmath, amssymb}
\usepackage{booktabs}
\usepackage{xcolor}
\usepackage{listings}
\usepackage[skins, breakable]{tcolorbox}

% --- Book colour palette (kept in sync with book/preamble.tex) ---
\definecolor{linkblue}{RGB}{14, 75, 140}
\definecolor{deepdivebg}{RGB}{235, 244, 255}
\definecolor{narrativebg}{RGB}{255, 252, 240}
\definecolor{narrativerule}{RGB}{180, 130, 40}
\definecolor{steelblue}{RGB}{70, 130, 180}
\definecolor{forestgreen}{RGB}{34, 139, 34}
\definecolor{crimson}{RGB}{220, 20, 60}
\definecolor{darkorange}{RGB}{255, 140, 0}
\definecolor{codebg}{RGB}{248, 248, 248}
\definecolor{coderule}{RGB}{210, 210, 210}
\definecolor{keywordblue}{RGB}{0, 100, 180}
\definecolor{commentgray}{RGB}{120, 120, 120}
\definecolor{stringgreen}{RGB}{30, 130, 30}

% Tie the Madrid structural colour to the book's link blue.
\setbeamercolor{structure}{fg=linkblue}
\setbeamercolor{title}{fg=white}
\setbeamercolor{frametitle}{fg=white}

% --- "the bigger picture" : narrative / intuition framing -----------
\newtcolorbox{bigpicture}{%
  enhanced, breakable, boxsep=2pt,
  colback  = narrativebg,
  colframe = narrativerule,
  boxrule  = 0pt, toprule = 1.4pt, bottomrule = 0.4pt,
  leftrule = 0pt, rightrule = 0pt, arc = 0pt,
  left = 6pt, right = 6pt, top = 4pt, bottom = 5pt,
  fonttitle = \footnotesize\scshape\color{narrativerule},
  title = {the bigger picture}, attach title to upper = {\par\smallskip},
}

% --- "under the hood" : the mechanism / technical detail ------------
\newtcolorbox{underhood}{%
  enhanced, breakable, boxsep=2pt,
  colback  = deepdivebg,
  colframe = linkblue,
  boxrule  = 0pt, toprule = 1.4pt, bottomrule = 0.4pt,
  leftrule = 0pt, rightrule = 0pt, arc = 0pt,
  left = 6pt, right = 6pt, top = 4pt, bottom = 5pt,
  fonttitle = \footnotesize\scshape\color{linkblue},
  title = {under the hood}, attach title to upper = {\par\smallskip},
}

% --- Python code style ---------------------------------------------
\lstset{
  basicstyle    = \ttfamily\footnotesize,
  keywordstyle  = \color{keywordblue}\bfseries,
  commentstyle  = \color{commentgray}\itshape,
  stringstyle   = \color{stringgreen},
  showstringspaces = false,
  language      = Python,
  breaklines    = true,
  backgroundcolor = \color{codebg},
  frame         = single,
  rulecolor     = \color{coderule},
  numbers       = none,
  aboveskip     = 4pt, belowskip = 4pt,
}

% --- Appendix conventions ------------------------------------------
% \appendixstart{Title} : begins the "extra material" appendix. Frames
% after it are not counted toward the main deck's frame total, keeping
% the core talk lean while preserving the deeper / optional content.
\newcommand{\appendixstart}[1]{%
  \appendix
  \setbeamertemplate{frametitle continuation}{}%
  \begin{frame}[noframenumbering]
    \begin{center}
      {\Large\scshape\color{linkblue} Appendix}\\[4pt]
      {\large #1}\\[10pt]
      {\footnotesize Extra / optional material — beyond the core lecture.}
    \end{center}
  \end{frame}
}
% Use \begin{frame}[noframenumbering]{...} for each appendix frame.

%   \title{...}\subtitle{...}\author{...}\institute{...}\date{\today}
%   \begin{document} ... \end{document}
% =====================================================================

\usetheme{Madrid}
\usecolortheme{whale}
\setbeamertemplate{navigation symbols}{}
\setbeamertemplate{caption}[numbered]
\setbeamercovered{transparent}

% --- Density controls: keep dense, equation-heavy frames on one slide ---
% Slightly smaller body text + tighter lists/blocks. Frame titles and the
% Madrid chrome keep their normal size, so the deck still reads as Madrid,
% just with more vertical room for content.
\setbeamerfont{normal text}{size=\small}
\setbeamertemplate{itemize/enumerate body begin}{\small}
\setbeamertemplate{itemize/enumerate subbody begin}{\footnotesize}
% Tighter block spacing.
\setbeamertemplate{block begin}{%
  \par\vskip2pt%
  \begin{beamercolorbox}[colsep*=.75ex,leftskip=1ex,rightskip=1ex]{block title}%
    \usebeamerfont*{block title}\insertblocktitle%
  \end{beamercolorbox}%
  {\parskip0pt\vskip-1pt}%
  \begin{beamercolorbox}[colsep*=.75ex,vmode,leftskip=1ex,rightskip=1ex]{block body}%
  \vskip1pt}
\setbeamertemplate{block end}{\end{beamercolorbox}\vskip2pt}

\usepackage{amsmath, amssymb}
\usepackage{booktabs}
\usepackage{xcolor}
\usepackage{listings}
\usepackage[skins, breakable]{tcolorbox}

% --- Book colour palette (kept in sync with book/preamble.tex) ---
\definecolor{linkblue}{RGB}{14, 75, 140}
\definecolor{deepdivebg}{RGB}{235, 244, 255}
\definecolor{narrativebg}{RGB}{255, 252, 240}
\definecolor{narrativerule}{RGB}{180, 130, 40}
\definecolor{steelblue}{RGB}{70, 130, 180}
\definecolor{forestgreen}{RGB}{34, 139, 34}
\definecolor{crimson}{RGB}{220, 20, 60}
\definecolor{darkorange}{RGB}{255, 140, 0}
\definecolor{codebg}{RGB}{248, 248, 248}
\definecolor{coderule}{RGB}{210, 210, 210}
\definecolor{keywordblue}{RGB}{0, 100, 180}
\definecolor{commentgray}{RGB}{120, 120, 120}
\definecolor{stringgreen}{RGB}{30, 130, 30}

% Tie the Madrid structural colour to the book's link blue.
\setbeamercolor{structure}{fg=linkblue}
\setbeamercolor{title}{fg=white}
\setbeamercolor{frametitle}{fg=white}

% --- "the bigger picture" : narrative / intuition framing -----------
\newtcolorbox{bigpicture}{%
  enhanced, breakable, boxsep=2pt,
  colback  = narrativebg,
  colframe = narrativerule,
  boxrule  = 0pt, toprule = 1.4pt, bottomrule = 0.4pt,
  leftrule = 0pt, rightrule = 0pt, arc = 0pt,
  left = 6pt, right = 6pt, top = 4pt, bottom = 5pt,
  fonttitle = \footnotesize\scshape\color{narrativerule},
  title = {the bigger picture}, attach title to upper = {\par\smallskip},
}

% --- "under the hood" : the mechanism / technical detail ------------
\newtcolorbox{underhood}{%
  enhanced, breakable, boxsep=2pt,
  colback  = deepdivebg,
  colframe = linkblue,
  boxrule  = 0pt, toprule = 1.4pt, bottomrule = 0.4pt,
  leftrule = 0pt, rightrule = 0pt, arc = 0pt,
  left = 6pt, right = 6pt, top = 4pt, bottom = 5pt,
  fonttitle = \footnotesize\scshape\color{linkblue},
  title = {under the hood}, attach title to upper = {\par\smallskip},
}

% --- Python code style ---------------------------------------------
\lstset{
  basicstyle    = \ttfamily\footnotesize,
  keywordstyle  = \color{keywordblue}\bfseries,
  commentstyle  = \color{commentgray}\itshape,
  stringstyle   = \color{stringgreen},
  showstringspaces = false,
  language      = Python,
  breaklines    = true,
  backgroundcolor = \color{codebg},
  frame         = single,
  rulecolor     = \color{coderule},
  numbers       = none,
  aboveskip     = 4pt, belowskip = 4pt,
}

% --- Appendix conventions ------------------------------------------
% \appendixstart{Title} : begins the "extra material" appendix. Frames
% after it are not counted toward the main deck's frame total, keeping
% the core talk lean while preserving the deeper / optional content.
\newcommand{\appendixstart}[1]{%
  \appendix
  \setbeamertemplate{frametitle continuation}{}%
  \begin{frame}[noframenumbering]
    \begin{center}
      {\Large\scshape\color{linkblue} Appendix}\\[4pt]
      {\large #1}\\[10pt]
      {\footnotesize Extra / optional material — beyond the core lecture.}
    \end{center}
  \end{frame}
}
% Use \begin{frame}[noframenumbering]{...} for each appendix frame.


\title{Practical Session 15: Privacy \& Local Models}
\subtitle{Hands-On: DP Budgets, NER De-identification, and a Membership-Inference Audit}
\author{Juan F.\ Imbet}
\institute{Paris Dauphine -- PSL University}
\date{\today}

\begin{document}

\begin{frame}\titlepage\end{frame}

% =====================================================================
\begin{frame}{Session Overview}
\textbf{Duration:} 1 hour (50 min problems + 10 min debrief)

\medskip
\textbf{Timed agenda:}
\begin{enumerate}
  \item \textbf{Recap} (8 min): DP definition, Gaussian mechanism, advanced composition,
        membership-inference AUC.
  \item \textbf{Problem 1} (12 min): Size the DP-SGD noise $\sigma$ and spend a
        privacy budget across multiple fine-tunes.
  \item \textbf{Problem 2} (12 min): Build a financial NER de-identification mask by
        hand; reason about $k$-anonymity vs.\ contextual re-identification.
  \item \textbf{Case Study} (15 min): Run the membership-inference audit notebook in
        \texttt{code/notebooks/15-privacy-local-models/}; read the AUC.
  \item \textbf{Discussion} (10 min): on-prem vs.\ private cloud vs.\ anonymise-and-send.
\end{enumerate}

\medskip
\begin{block}{Goal}
Leave able to (a) cost a privacy budget, (b) hand-build a de-identification mask, and
(c) interpret an empirical leakage metric.
\end{block}
\end{frame}

% =====================================================================
\begin{frame}{Recap: The Formulas You Will Need}
\begin{block}{Gaussian mechanism (sets the noise scale)}
\[
\sigma = \frac{\Delta_2 f \cdot \sqrt{2\ln(1.25/\delta)}}{\varepsilon},
\qquad \mathcal{M}(D) = f(D) + \mathcal{N}(0,\sigma^2 I)
\]
\end{block}

\begin{block}{Advanced composition over $k$ steps}
\[
\varepsilon' \approx \varepsilon\sqrt{2k\ln(1/\delta')}
\]
\end{block}

\textbf{DP-SGD clip}: $\hat{g}_i = g_i / \max(1, \|g_i\|_2 / C)$ --- bounds per-example
sensitivity to $C$.

\medskip
\textbf{Membership-inference AUC}: $0.5$ = no leakage, $1.0$ = total failure;
$<0.6$ commonly acceptable.
\end{frame}

% =====================================================================
\begin{frame}{Recap: De-identification Toolkit}
\begin{columns}[T]
\begin{column}{0.5\textwidth}
\textbf{Two-layer NER pipeline}
\begin{itemize}
  \item Learned model: PERSON, ORG, LOCATION, MONEY, DATE
  \item Rules / regex: IBAN, ISIN, LEI, tax ID, SWIFT
\end{itemize}
\medskip
\textbf{Three strategies}
\begin{itemize}
  \item Masking: \texttt{[PERSON\_001]}
  \item Generalisation: 75008 $\to$ Paris
  \item Suppression: remove span
\end{itemize}
\end{column}
\begin{column}{0.5\textwidth}
\begin{block}{$k$-anonymity}
Output could plausibly come from $\geq k$ individuals. Masking gives large $k$ ---
\emph{but} context can still re-identify.
\end{block}
\medskip
\textbf{Pseudonymisation} keeps text natural but the mapping is a key: GDPR
Recital 26 $\Rightarrow$ protected, not anonymised.
\end{column}
\end{columns}
\end{frame}

% =====================================================================
\begin{frame}{Problem 1: Spending a Privacy Budget}
\textbf{Scenario.} A bank fine-tunes a 7B model on credit-decision narratives with
DP-SGD. Per-example gradients are clipped to $C = 1.0$. Target per-run guarantee:
$\varepsilon_{\text{run}} = 0.8$, $\delta = 10^{-5}$.

\medskip
\textbf{Task A.} The aggregated clipped gradient has $\ell_2$-sensitivity
$\Delta_2 f = C = 1.0$. Compute the noise scale $\sigma$ for one mechanism
application using the Gaussian mechanism. Use $\ln(1.25/10^{-5}) \approx 11.7$.

\medskip
\textbf{Task B.} The team plans $k = 4$ independent fine-tuning runs over the year,
each $(\varepsilon_{\text{run}}, \delta)$-DP. Using advanced composition with
$\delta' = 10^{-5}$ (so $\ln(1/\delta') \approx 11.5$), what is the total budget
$\varepsilon'$?

\medskip
\textbf{Task C.} The institution's annual privacy budget ceiling is
$\varepsilon_{\max} = 5$. Are the four runs within budget? How many such runs can it
afford before the ceiling is breached?
\end{frame}

% =====================================================================
\begin{frame}{Problem 2: Build a De-identification Mask}
\textbf{Source sentence (relationship-manager call note):}

\medskip
\footnotesize
\emph{``Spoke with Claire Moreau (LEI 969500TEST00CLAIRE12) about increasing the
EUR 1,250,000 revolving facility; she will wire from IBAN
DE89 3704 0044 0532 0130 00 once the board of Helvetia Trading SA approves on
12 March 2025.''}

\medskip
\normalsize
\textbf{Task A.} List every sensitive span and assign each a type
(PERSON, ORG, MONEY, DATE, IBAN, LEI). Mark which need the \emph{learned} model and
which the \emph{rule} layer.

\medskip
\textbf{Task B.} Produce the fully \emph{masked} version with typed placeholders
(\texttt{[PERSON\_001]}, etc.), keeping co-references consistent.

\medskip
\textbf{Task C.} Suppose the portfolio contains exactly one Swiss trading firm with a
revolving facility near EUR 1.25m maturing in Q1 2025. Does masking achieve
$k$-anonymity here? Which strategy (generalisation / suppression) would you add, and
to which span?
\end{frame}

% =====================================================================
\begin{frame}{Case Study: Membership-Inference Audit}
\textbf{Goal.} Quantify how much a fine-tuned model leaks about its training set ---
the empirical check the chapter insists you run \emph{after every fine-tune}.

\medskip
\textbf{Notebook:} \texttt{code/notebooks/15-privacy-local-models/exercises.ipynb}\\
\textbf{Figure generator:} \texttt{gen\_dp\_privacy\_utility.py} (privacy--utility curve).

\medskip
\textbf{The procedure:}
\begin{enumerate}
  \item Hold out a random subset of fine-tuning documents as the \emph{non-member}
        (negative) class.
  \item Score every member and non-member with the likelihood-ratio attack
        (target vs.\ reference model).
  \item Sweep the decision threshold; plot the ROC; report the \textbf{AUC}.
\end{enumerate}

\medskip
\begin{alertblock}{Acceptance gate}
AUC $\approx 0.5$ $\Rightarrow$ attack no better than chance (good). AUC $> 0.6$
$\Rightarrow$ leakage: add DP-SGD, reduce epochs, or exclude the most sensitive docs.
\end{alertblock}
\end{frame}

% =====================================================================
\begin{frame}[fragile]{Case Study: The Audit Code (1/2 --- the attack)}
\begin{lstlisting}
import numpy as np
from sklearn.metrics import roc_auc_score

# Per-document loss under target (fine-tuned) and reference models.
# Lower target loss on a doc => evidence it was a training member.
def membership_scores(target_loss, ref_loss):
    # Likelihood-ratio signal: high when target fits doc better than ref.
    return ref_loss - target_loss          # larger => more "member-like"

def audit_auc(target_loss, ref_loss, is_member):
    scores = membership_scores(target_loss, ref_loss)
    return roc_auc_score(is_member, scores)
\end{lstlisting}
\end{frame}

% =====================================================================
\begin{frame}[fragile]{Case Study: The Audit Code (2/2 --- the toy run)}
\begin{lstlisting}
# Toy illustration: members fit noticeably better under the target model.
rng = np.random.default_rng(0)
ref     = rng.normal(2.0, 0.4, size=400)
member  = ref[:200] - rng.normal(0.6, 0.2, size=200)   # over-fit members
nonmemb = ref[200:] - rng.normal(0.0, 0.2, size=200)   # no advantage
target  = np.concatenate([member, nonmemb])
is_member = np.array([1]*200 + [0]*200)

auc = audit_auc(target, ref, is_member)
print(f"membership-inference AUC = {auc:.3f}")
# Interpret: <0.6 acceptable; ~0.5 ideal; ->1.0 means severe leakage.
\end{lstlisting}
\end{frame}

% =====================================================================
\begin{frame}{Case Study: Diagnostic Questions}
Run the audit and the privacy--utility generator, then answer:

\medskip
\textbf{Q1.} What AUC does the toy run report? Does it indicate acceptable privacy?
Which array (\texttt{member} vs.\ \texttt{nonmemb}) drives the gap, and why?

\medskip
\textbf{Q2.} Re-run \texttt{member} with the over-fit offset reduced from $0.6$ to
$0.0$. What happens to the AUC, and what does that tell you about the link between
over-fitting and membership leakage?

\medskip
\textbf{Q3.} In \texttt{gen\_dp\_privacy\_utility.py}, the curve is
$\text{acc}(\varepsilon)=0.90-(0.90-0.55)e^{-0.45\varepsilon}$. Compute the modelled
accuracy at $\varepsilon=1$ and $\varepsilon=10$. How large is the utility gap to the
non-private baseline (0.90) in each case?

\medskip
\textbf{Q4 (link the two).} You must keep MI-AUC $<0.6$ \emph{and} accuracy $>0.80$.
Reading both results, is a single DP fine-tune at $\varepsilon=1$ likely to satisfy
both? What lever (clip norm $C$, epochs, LoRA) would you pull first?
\end{frame}

% =====================================================================
\begin{frame}{Discussion: Choosing the Deployment Pattern}
For each scenario, pick \textbf{on-premise}, \textbf{private cloud}, or
\textbf{anonymise-and-send to a cloud API} --- and justify against regulation,
sensitivity, latency, and cost.

\medskip
\begin{enumerate}
  \item A relationship-manager assistant that routinely sees client names, account
        numbers, and transaction histories. Latency target: interactive.
  \item Summarising \emph{public} ECB rate decisions for an internal newsletter ---
        no client data, high volume, cost-sensitive.
  \item Joint fraud-detection model across three banks that legally cannot share raw
        transaction data with one another.
  \item Drafting a bond prospectus ``Risk Factors'' section from a structured issuer
        data dump and five precedent deals (internal, confidential, not personal data).
\end{enumerate}

\medskip
\begin{block}{Frame your answer around}
External boundary $\cdot$ data-residency $\cdot$ record-keeping (MiFID II) $\cdot$
DORA third-party terms $\cdot$ memorisation risk $\cdot$ latency $\cdot$ cost.
\end{block}
\end{frame}

% =====================================================================
\begin{frame}{Solution: Problem 1 --- Privacy Budget}
\textbf{Task A --- noise scale} ($\Delta_2 f = 1$, $\varepsilon_{\text{run}}=0.8$,
$\ln(1.25/10^{-5})\approx 11.7$):
\[
\sigma = \frac{1 \cdot \sqrt{2 \times 11.7}}{0.8}
       = \frac{\sqrt{23.4}}{0.8} \approx \frac{4.84}{0.8} \approx 6.05
\]

\textbf{Task B --- four runs, advanced composition}
($k=4$, $\varepsilon=0.8$, $\ln(1/\delta')\approx 11.5$):
\[
\varepsilon' \approx \varepsilon\sqrt{2k\ln(1/\delta')}
= 0.8\sqrt{2\times4\times11.5}
= 0.8\sqrt{92} \approx 0.8 \times 9.59 \approx 7.67
\]

\textbf{Task C --- against the ceiling} $\varepsilon_{\max}=5$:
\[
7.67 > 5 \;\Rightarrow\; \textbf{four runs exceed budget.}
\]
Largest $k$ with $0.8\sqrt{2k\cdot 11.5}\le 5$: $\sqrt{23k}\le 6.25
\Rightarrow 23k \le 39.1 \Rightarrow k \le 1.7$ $\Rightarrow$ \textbf{only 1 run}
fits under advanced composition at this $\varepsilon_{\text{run}}$. Lesson: budget is
scarce --- lower $\varepsilon_{\text{run}}$ per run or fine-tune less often.
\end{frame}

% =====================================================================
\begin{frame}{Solution: Problem 2 --- De-identification Mask}
\textbf{Task A --- spans and types:}
\footnotesize
\begin{tabular}{lll}
\toprule
Span & Type & Layer \\
\midrule
Claire Moreau & PERSON & Learned \\
969500TEST00CLAIRE12 & LEI & Rule \\
EUR 1,250,000 & MONEY & Learned \\
DE89 3704 0044 0532 0130 00 & IBAN & Rule \\
Helvetia Trading SA & ORG & Learned \\
12 March 2025 & DATE & Learned \\
\bottomrule
\end{tabular}

\medskip
\normalsize
\textbf{Task B --- masked output:}
\footnotesize
\emph{``Spoke with [PERSON\_001] (LEI [LEI\_001]) about increasing the [MONEY\_001]
revolving facility; she will wire from IBAN [IBAN\_001] once the board of [ORG\_001]
approves on [DATE\_001].''}

\medskip
\normalsize
\textbf{Task C --- re-identification.} Masking does \emph{not} achieve $k$-anonymity:
``one Swiss trading firm, $\sim$EUR 1.25m revolver, Q1 2025'' is a near-unique
fingerprint. Add \textbf{generalisation} --- MONEY $\to$ ``EUR 1--2m'', DATE $\to$
``Q1 2025'' --- and consider \textbf{suppressing} the ORG nationality cue. Goal:
posterior over individuals stays diffuse.
\end{frame}

% =====================================================================
\appendixstart{Stretch Problems}
% =====================================================================

\begin{frame}[noframenumbering]{Appendix --- Stretch Problem A: Gradient Inversion}
Federated learning shares only gradient updates, never raw data --- yet
\textbf{gradient inversion} can reconstruct inputs from a single update.

\medskip
\textbf{A1.} Explain why a FedAvg update $\Delta_k^t$ from a client training on one
sensitive batch can leak that batch, and why \emph{larger} local batch size and
\emph{more} local steps $E$ both reduce the leakage.

\medskip
\textbf{A2.} You add \emph{local} DP-SGD before transmitting $\Delta_k^t$. State the
two new costs this imposes (relative to plain FedAvg) and why they are unavoidable.

\medskip
\textbf{A3.} A cross-border bank runs one federated client per jurisdiction. Which
regulation does this architecture specifically satisfy that pure centralised training
would violate, and why does federation alone \emph{not} discharge GDPR Art.\ 17?
\end{frame}

\begin{frame}[noframenumbering]{Appendix --- Stretch Problem B: Quantisation Audit}
A risk team must sign off a 4-bit (AWQ) quantised, DP-fine-tuned model for on-prem
deployment.

\medskip
\textbf{B1.} The team worries quantisation might \emph{change} the privacy guarantee.
Argue, using the post-processing immunity property of DP, why quantising a
DP-trained model does \emph{not} weaken its $(\varepsilon,\delta)$ guarantee.

\medskip
\textbf{B2.} AWQ needs a calibration set. Why is using a small \emph{public} corpus
(not client data) both sufficient for quality and necessary for privacy?

\medskip
\textbf{B3.} 7B at 4-bit $\approx 3.5$\,GB. A branch workstation has a single 8\,GB
GPU. Estimate whether a 13B 4-bit model ($\approx 6.5$\,GB) also fits, and name one
operational reason you might still prefer the 7B at the branch.
\end{frame}

\end{document}
